\documentclass[11pt,a4paper]{article}
\usepackage[utf8]{inputenc}
\usepackage[T1]{fontenc}
\usepackage[english]{babel}
\usepackage[margin=1.3cm,top=0.9cm,bottom=0.8cm]{geometry}
\usepackage{hyperref}
\usepackage{titlesec}
\usepackage{enumitem}
\setlist{topsep=2pt}
\usepackage{xcolor}
\usepackage{amssymb}

% Colours and section formatting
\definecolor{primary}{RGB}{30, 80, 130}
\titleformat{\section}{\Large\bfseries\color{primary}}{}{0em}{}[\titlerule]
\titlespacing{\section}{0pt}{0.9ex}{0.5ex}

\begin{document}
	\pagestyle{empty}
	
	% Header
	\begin{center}
		{\Huge \textbf{Yanis Bérard}}\\[0.2cm]
		{\Large Student at the École Normale Supérieure de Rennes -- Mathematics}\\[0.2cm]
		\href{mailto:yanis.berard@ens-rennes.fr}{yanis.berard@ens-rennes.fr} ~|~ +33 (0)7 68 83 62 00 ~|~ \href{https://yanisberard-sys.github.io/}{yanisberard-sys.github.io}
	\end{center}
	
	\vspace{0cm}
	
	\section*{Education}
	\begin{itemize}[leftmargin=1em, label={}]
		\item \textbf{ENS Rennes} \hfill \textit{2025 -- Present} \\
		Student (\emph{élève normalien}), Pure Mathematics track (\emph{Magistère de Mathématiques}).
		\vspace{0.02cm}
		
		\item \textbf{Université de Rennes} \hfill \textit{2025 -- Present} \\
		Master's degree in Mathematics, first year (M1, 2026 -- Present); Bachelor's degree in Mathematics, third year (Licence 3, 2025 -- 2026).
		\vspace{0.02cm}
		
		\item \textbf{Lycée du Parc, Lyon} \hfill \textit{2024 -- 2025} \\
		CPGE MP*1 (two-year undergraduate preparatory classes for the Grandes Écoles).
		\vspace{0.02cm}
		
		\item \textbf{Lycée La Martinière Monplaisir, Lyon} \hfill \textit{2022 -- 2024} \\
		CPGE MPSI $\rightarrow$ MP*.
		\vspace{0.02cm}
		
		\item \textbf{Lycée Honoré d'Urfé, Saint-Étienne} \hfill \textit{2022} \\
		Baccalauréat, Mathematics and Physics majors, with highest honours (\emph{mention Très Bien}).
	\end{itemize}
	
	\section*{Research Experience}
	\begin{itemize}[leftmargin=1em, label={}]
		\item \textbf{Research internship, Institut Camille Jordan, Lyon} \hfill \textit{May -- June 2026} \\
		Supervised by Jean-Yves Welschinger. Report: \textit{Local Holomorphy and Global Algebraicity: The Road to GAGA} (compact Riemann surfaces, covering theory and the Riemann--Hurwitz formula, sheaf cohomology, Riemann--Roch theorem, Serre duality, projective embedding, Chow's theorem, GAGA for curves).
		\vspace{0.02cm}
		
		\item \textbf{GATE Laboratory, Saint-Étienne} \hfill \textit{2023} \\
		Study of voting theory and numerical simulations around Arrow's impossibility theorem.
	\end{itemize}
	
	\section*{Academic Projects}
	\begin{itemize}[leftmargin=1.5em, label=\textbullet, itemsep=2pt]
		\item \textbf{Expository paper (2026):} \textit{An elementary proof of Bézout's theorem on algebraic curves}. Study of intersections of plane projective curves, local intersection multiplicities, explicit examples and geometric interpretation. Three-hour presentation to fellow students.
		\item \textbf{Expository paper:} \textit{The Kakeya conjecture} over finite and local fields, and over the reals.
		\item \textbf{Group theory notes:} simple groups of order $< 60$; duality for finite abelian groups.
		\item \textbf{MATh.en.JEANS (2022):} recovering a triangle from its special points.
	\end{itemize}

	\section*{TIPE (supervised research projects in CPGE)}
	\begin{itemize}[leftmargin=1.5em, label=\textbullet, itemsep=2pt]
		\item \textbf{Elections: the rules of the democratic game} (2022--23): Arrow's theorem, computational social choice.
		\item \textbf{Heads or tails: paradoxes and random walks} (2023--24): waiting times, Pólya's theorem.
		\item \textbf{Lattices of finite groups and applications} (2024--25): a generation game on finite groups.
	\end{itemize}
	
	\section*{Teaching, Outreach \& Other Experience}
	\begin{itemize}[leftmargin=1em, label={}]
		\item \textbf{Private tutor in Mathematics and Physics} \hfill \textit{2021 -- Present} \\
		Regular tutoring of students from middle school to preparatory classes (CPGE), including preparation for the competitive entrance exams of the Grandes Écoles.
		\vspace{0.02cm}
		
		\item \textbf{Outreach:} Organisation of the \emph{Rencontres des Jeunes Mathématiciennes et Informaticiennes} (outreach event for young women in mathematics and computer science), design of a workshop.
		\vspace{0.02cm}
		
		\item \textbf{Mechanical worker} -- \textit{NTD France}; \textbf{Market gardening} -- \textit{Ferme Le Bénitier aux Oiseaux}
	\end{itemize}
	
	\section*{Skills}
	\begin{itemize}[leftmargin=1.5em, label=\textbullet, itemsep=2pt]
		\item \textbf{Languages:} French (native), English (C1), Portuguese (B2), Spanish (B1), Italian (B1).
		\item \textbf{Computer skills:} \LaTeX, Python, OCaml.
		\item \textbf{Interests:} Climbing (Loire junior champion), skiing, running, cycling, chess.
	\end{itemize}
	
\end{document}
